\documentclass{beamer}

\usepackage{beamerthemeSingapore}
\usepackage[brazil]{babel}
\usepackage[T1]{fontenc}
\usepackage[utf8]{inputenc}
\usepackage{graphicx}
\usepackage{url}
\usepackage{xcolor}
\usepackage{booktabs}
\usepackage{ragged2e}
\usepackage{etoolbox}
\usepackage{amsmath}
\usepackage{minted}
\usepackage{tikz}
\usetikzlibrary{positioning, arrows.meta, shapes.geometric}

\newcommand{\tcb}[1]{\textcolor{blue}{#1}}
\newcommand{\tcr}[1]{\textcolor{red}{#1}}
\newcommand{\tcg}[1]{\textcolor{green!50!black}{#1}}

\tikzset{
  treenode/.style = {align=center, inner sep=0pt, text centered,
    font=\sffamily\small},
  no/.style = {treenode, circle, black, draw=black, fill=white,
    text width=1.7em, line width=0.6pt},
  lv/.style = {draw=red!80!black, line width=2.2pt},
  lp/.style = {draw=black, line width=0.6pt},
  arn_t/.style = {treenode, regular polygon, regular polygon sides=3,
    draw=black, fill=black!8, minimum size=1.1cm, inner sep=0pt, line width=0.6pt,
    font=\sffamily\footnotesize},
  bnode/.style = {draw=black, fill=black!5, rounded corners=1pt,
    font=\sffamily\footnotesize, inner sep=3pt},
  tnode/.style = {treenode, circle, draw=black, fill=white,
    minimum size=1.5em, inner sep=1pt, font=\sffamily\footnotesize},
  tfim/.style = {treenode, circle, draw=black, fill=green!25,
    minimum size=1.5em, inner sep=1pt, font=\sffamily\footnotesize}
}

\setminted{
  fontsize=\footnotesize,
  breaklines=true,
  breakanywhere=true,
  frame=lines,
  framesep=2mm,
  baselinestretch=1.05,
  tabsize=2
}

\apptocmd{\frame}{}{\justifying}{}

\setbeamertemplate{navigation symbols}{}

\title{09 - Árvores rubro-negras, árvores B e tries}
\author{Prof. Alex Kutzke}
\date{Setembro de 2026}

\begin{document}

% ==========================================
\begin{frame}
  \begin{center}
    \large{Setor de Educação Profissional e Tecnológica}\\
    Tecnologia em Análise e Desenvolvimento de Sistemas\\
    \vspace{1cm}
    \small{DS143 - Estruturas de Dados II}
  \end{center}
  \titlepage
\end{frame}

% ==========================================
\begin{frame}{Objetivos da aula}
  Ao final desta aula você deve ser capaz de:
  \begin{enumerate}
    \item Enunciar as propriedades de uma árvore rubro-negra e verificar se uma
          árvore colorida dada as satisfaz;
    \item Explicar por que essas propriedades limitam a altura a
          $2\log_2(n+1)$;
    \item Descrever a inserção em termos de recoloração e rotação;
    \item Comparar AVL e rubro-negra, e escolher entre as duas a partir do
          perfil de uso;
    \item Explicar por que a árvore binária é a estrutura errada em disco, e
          como a árvore B resolve o problema;
    \item Executar à mão a inserção com divisão de nó em uma árvore B;
    \item Descrever a trie e justificar por que o custo de busca nela não
          depende do número de chaves.
  \end{enumerate}

  \vspace{0.2cm}
  \footnotesize Este conteúdo cai na Prova 2.
\end{frame}

% ==========================================
\begin{frame}{Roteiro}
  \tableofcontents
\end{frame}

% ==========================================
\section{Rubro-negras}

\begin{frame}[fragile]{Demonstração: o índice de um banco de dados}
\begin{minted}{sql}
CREATE TABLE aluno(grr INTEGER, nome TEXT);
WITH RECURSIVE seq(n) AS (
  SELECT 1 UNION ALL SELECT n + 1 FROM seq
  WHERE n < 2000000
)
INSERT INTO aluno(grr, nome)
  SELECT n, 'aluno ' || n FROM seq;

EXPLAIN QUERY PLAN
  SELECT nome FROM aluno WHERE grr = 1999999;

CREATE INDEX idx_grr ON aluno(grr);

EXPLAIN QUERY PLAN
  SELECT nome FROM aluno WHERE grr = 1999999;
\end{minted}

  \vspace{0.2cm}
  Dois milhões de linhas, a mesma consulta duas vezes.
\end{frame}

\begin{frame}[fragile]{O plano de execução muda}
\begin{minted}{text}
QUERY PLAN
`--SCAN aluno
aluno 1999999
Run Time: real 0.069

QUERY PLAN
`--SEARCH aluno USING INDEX idx_grr (grr=?)
aluno 1999999
Run Time: real 0.000

4096            <- PRAGMA page_size
\end{minted}

  \vspace{0.2cm}
  De 69 ms para menos de 1 ms. O índice é uma \textbf{árvore B}, e 4096 é o
  tamanho em bytes de cada nó dela. Voltamos a esta tela no fim da aula.
\end{frame}

\begin{frame}{Três estruturas, a partir da ABB balanceada}
  \begin{block}{Árvore rubro-negra}
    Mantém a garantia de $O(\log n)$ com um critério de balanceamento mais
    frouxo que o da AVL. É a árvore de busca das bibliotecas de linguagem.
  \end{block}

  \begin{block}{Árvore B}
    Abandona o formato binário para funcionar em disco. É a estrutura de índice
    de bancos de dados e de sistemas de arquivos.
  \end{block}

  \begin{block}{Trie}
    Abandona a comparação entre chaves inteiras e usa a estrutura interna da
    chave, letra por letra.
  \end{block}
\end{frame}

\begin{frame}{O que o critério da AVL cobra}
  Alturas das duas subárvores diferindo em no máximo 1, em todo nó.

  \vspace{0.4cm}
  \begin{itemize}
    \item um campo de altura por nó, atualizado em todo nó do caminho;
    \item a remoção pode rotacionar em cada nível até a raiz: até
          $O(\log n)$ rotações para tirar \tcr{uma} chave.
  \end{itemize}

  \vspace{0.4cm}
  A inserção é barata (no máximo uma rotação simples ou dupla). A remoção é que
  decide em uma estrutura que recebe alterações o tempo todo.

  \vspace{0.4cm}
  \textbf{Pergunta:} quanto desbalanceamento dá para admitir sem perder a altura
  logarítmica?
\end{frame}

\begin{frame}{As três propriedades}
  \begin{columns}
    \begin{column}{0.5\textwidth}
      \centering
      \begin{tikzpicture}[-, level distance=0.95cm,
          level 1/.style={sibling distance=2.6cm},
          level 2/.style={sibling distance=1.3cm},
          level 3/.style={sibling distance=0.8cm}]
        \node [no] {40}
        child { node [no] {30}
            child { node [no] {20}
                child { node [no] {10} edge from parent [lv] }
                child [missing]
              }
            child { node [no] {35} }
          }
        child { node [no] {50}
            child { node [no] {45} }
            child { node [no] {95}
                child { node [no] {90} edge from parent [lv] }
                child [missing]
              }
          };
      \end{tikzpicture}
    \end{column}
    \begin{column}{0.5\textwidth}
      \small
      Uma ABB em que cada nó tem uma cor e vale:
      \begin{enumerate}
        \item a raiz é preta;
        \item nenhum nó vermelho tem filho vermelho;
        \item todo caminho da raiz até uma subárvore vazia tem o mesmo número
              de nós pretos.
      \end{enumerate}

      \vspace{0.2cm}
      Esse número é a \textbf{altura preta}. Aqui, 3.

      \vspace{0.2cm}
      \footnotesize Desenho de Sedgewick: a cor do nó é a da \textbf{ligação}
      que vem do pai. Vermelhos: 10 e 90.
    \end{column}
  \end{columns}
\end{frame}

\begin{frame}{Por que a altura é logarítmica}
  Seja $b$ a altura preta.

  \vspace{0.3cm}
  \begin{itemize}
    \item o caminho mais curto só tem pretos: $b$ nós;
    \item o mais longo alterna preto e vermelho (propriedade 2): no máximo
          $2b$ nós.
  \end{itemize}

  \vspace{0.3cm}
  Nenhum caminho passa do dobro de outro. Com pelo menos $2^b - 1$ nós na
  árvore, vale $b \le \log_2(n+1)$ e

  \vspace{0.2cm}
  \begin{center}
    \fbox{$h \le 2\log_2(n+1)$}
  \end{center}

  \vspace{0.3cm}
  Contra $1{,}44\log_2(n+2)$ da AVL. O limite é pior; o preço de manter a
  árvore dentro dele é menor.
\end{frame}

\begin{frame}{As cores codificam uma árvore 2-3}
  \begin{columns}
    \begin{column}{0.42\textwidth}
      \centering
      \begin{tikzpicture}[-, level distance=1.1cm,
          level 1/.style={sibling distance=1.6cm}]
        \node [bnode] {20 $\mid$ 40}
        child { node [arn_t] {} }
        child { node [arn_t] {} }
        child { node [arn_t] {} };
      \end{tikzpicture}

      \vspace{0.3cm}
      \footnotesize Nó 2-3 com duas chaves e três filhos.
    \end{column}
    \begin{column}{0.42\textwidth}
      \centering
      \begin{tikzpicture}[-, level distance=1cm,
          level 1/.style={sibling distance=1.8cm},
          level 2/.style={sibling distance=1.3cm}]
        \node [no] {40}
        child { node [no] {20}
            child { node [arn_t] {} edge from parent [lp] }
            child { node [arn_t] {} edge from parent [lp] }
            edge from parent [lv]
          }
        child { node [arn_t] {} };
      \end{tikzpicture}

      \vspace{0.3cm}
      \footnotesize A ligação \tcr{vermelha} une as duas chaves do mesmo nó 2-3.
    \end{column}
  \end{columns}

  \vspace{0.4cm}
  \begin{itemize}
    \item ligação preta = descer um nível na árvore 2-3;
    \item altura preta = altura da árvore 2-3.
  \end{itemize}
\end{frame}

\begin{frame}{A variante left-leaning}
  Sedgewick (seção 3.3) acrescenta duas exigências às três propriedades:

  \vspace{0.3cm}
  \begin{itemize}
    \item toda ligação vermelha aponta para a \textbf{esquerda};
    \item nenhum nó tem duas ligações vermelhas.
  \end{itemize}

  \vspace{0.3cm}
  Com elas, cada nó 2-3 tem uma única representação binária, e a correspondência
  com a árvore 2-3 é exata. A inserção perde metade dos casos e cabe em três
  testes.

  \vspace{0.3cm}
  O Cormen (cap. 13) admite vermelho à direita e nó preto com os dois filhos
  vermelhos, que é um nó de três chaves: as cores dele codificam uma árvore
  \textbf{2-3-4}.
\end{frame}

% ==========================================
\section{Inserção}

\begin{frame}{O nó novo entra vermelho}
  \begin{columns}
    \begin{column}{0.5\textwidth}
      \begin{block}{Se entrasse preto}
        Acrescentaria um nó preto ao caminho que o contém e quebraria a
        propriedade 3, a mais cara de restabelecer.
      \end{block}
    \end{column}
    \begin{column}{0.5\textwidth}
      \begin{block}{Entrando vermelho}
        Só pode quebrar a propriedade 2, e apenas quando o pai também for
        vermelho.
      \end{block}
    \end{column}
  \end{columns}

  \vspace{0.5cm}
  Os consertos acontecem na volta da recursão, em cada nó do caminho da folha
  nova até a raiz. Ao final, a raiz é pintada de preto.
\end{frame}

\begin{frame}{Os três consertos}
  \begin{center}
    \small
    \begin{tabular}{p{4.6cm}ll}
      \toprule
      \textbf{Situação no nó} & \textbf{Conserto} & \textbf{Muda} \\
      \midrule
      filho direito vermelho e esquerdo preto & rotação à esquerda & ponteiros \\
      filho esquerdo vermelho com filho esquerdo vermelho & rotação à direita & ponteiros \\
      os dois filhos vermelhos & inversão de cores & só as cores \\
      \bottomrule
    \end{tabular}
  \end{center}

  \vspace{0.4cm}
  Nessa ordem, sempre. A inversão de cores empurra o vermelho um nível para
  cima: na árvore 2-3, é a divisão de um nó que ficou com três chaves.
\end{frame}

\begin{frame}{Inversão de cores}
  \begin{columns}
    \begin{column}{0.45\textwidth}
      \centering
      \begin{tikzpicture}[-, level distance=1cm,
          level 1/.style={sibling distance=1.6cm}]
        \node [no] (r) {40}
        child { node [no] {20} edge from parent [lv] }
        child { node [no] {50} edge from parent [lv] };
        \draw [lp] (r.north) -- ++(0,0.6);
      \end{tikzpicture}

      \vspace{0.2cm}
      \footnotesize antes
    \end{column}
    \begin{column}{0.45\textwidth}
      \centering
      \begin{tikzpicture}[-, level distance=1cm,
          level 1/.style={sibling distance=1.6cm}]
        \node [no] (r) {40}
        child { node [no] {20} }
        child { node [no] {50} };
        \draw [lv] (r.north) -- ++(0,0.6);
      \end{tikzpicture}

      \vspace{0.2cm}
      \footnotesize depois
    \end{column}
  \end{columns}

  \vspace{0.4cm}
  Nenhum ponteiro muda. As duas ligações vermelhas de baixo viram pretas, e a
  ligação que chega ao 40 vira vermelha: o vermelho sobe um nível para ser
  tratado pelo pai. A altura preta dos caminhos não muda.
\end{frame}

\begin{frame}[fragile]{Rotação com cor}
\begin{minted}{go}
func rotacionaEsquerda(a *No) *No {
	b := a.Dir
	a.Dir = b.Esq
	b.Esq = a
	b.Cor = a.Cor
	a.Cor = Vermelho
	return b
}
\end{minted}

  \vspace{0.2cm}
  As três primeiras linhas são a rotação da AVL. As duas seguintes são o que a
  cor acrescenta: quem sobe leva a cor do nó que estava ali, e quem desce fica
  vermelho.
\end{frame}

\begin{frame}[fragile]{A inserção inteira}
\begin{minted}[fontsize=\scriptsize]{go}
func insere(a *No, v int) *No {
	if a == nil {
		return &No{Info: v, Cor: Vermelho}
	}
	if v < a.Info {
		a.Esq = insere(a.Esq, v)
	} else if v > a.Info {
		a.Dir = insere(a.Dir, v)
	}
	if vermelho(a.Dir) && !vermelho(a.Esq) {
		a = rotacionaEsquerda(a)
	}
	if vermelho(a.Esq) && vermelho(a.Esq.Esq) {
		a = rotacionaDireita(a)
	}
	if vermelho(a.Esq) && vermelho(a.Dir) {
		inverteCores(a)
	}
	return a
}
\end{minted}
\end{frame}

\begin{frame}{Inserindo 10, 20 e 30 em ordem}
  \begin{columns}[t]
    \begin{column}{0.3\textwidth}
      \centering
      \begin{tikzpicture}[-, level distance=0.9cm,
          level 1/.style={sibling distance=1.2cm}]
        \node [no] {10}
        child [missing]
        child { node [no] {20} edge from parent [lv] };
      \end{tikzpicture}

      \vspace{0.2cm}
      \footnotesize vermelho à direita
    \end{column}
    \begin{column}{0.3\textwidth}
      \centering
      \begin{tikzpicture}[-, level distance=0.9cm,
          level 1/.style={sibling distance=1.2cm}]
        \node [no] {20}
        child { node [no] {10} edge from parent [lv] }
        child [missing];
      \end{tikzpicture}

      \vspace{0.2cm}
      \footnotesize rotação à esquerda
    \end{column}
    \begin{column}{0.35\textwidth}
      \centering
      \begin{tikzpicture}[-, level distance=0.9cm,
          level 1/.style={sibling distance=1.3cm}]
        \node [no] {20}
        child { node [no] {10} }
        child { node [no] {30} };
      \end{tikzpicture}

      \vspace{0.2cm}
      \footnotesize insere 30, inversão de cores
    \end{column}
  \end{columns}

  \vspace{0.5cm}
  A entrada em ordem crescente, que degenera a ABB em uma lista, produz aqui uma
  árvore cheia. O 20 virou raiz sem ter sido inserido primeiro.
\end{frame}

\begin{frame}[fragile]{A árvore de exemplo, agora colorida}
\begin{minted}{text}
$ go run rn.go
Inserindo [50 30 90 20 40 95 10 35 45]

        95 (p)
            90 (v)
    50 (p)
        45 (p)
40 (p)
        35 (p)
    30 (p)
        20 (p)
            10 (v)

altura: 3   altura preta: 3   é rubro-negra: true
\end{minted}

  \vspace{0.1cm}
  \texttt{(v)}: ligação vermelha vinda do pai. É a árvore do slide das
  propriedades, e a raiz é o 40, e não o 50 que entrou primeiro.
\end{frame}

% ==========================================
\section{AVL e rubro-negra}

\begin{frame}[fragile]{Medindo as duas: altura e consertos}
\begin{minted}{text}
Inserções em ordem crescente
      n | alt. ABB | alt. AVL | alt. RN | rot. AVL | rot. RN
--------+----------+----------+---------+----------+--------
   1000 |      999 |        9 |       9 |      990 |     991
  10000 |     9999 |       13 |      13 |     9986 |    9987
 100000 |    99999 |       16 |      16 |    99983 |   99984

Inserções em ordem aleatória
      n | alt. ABB | alt. AVL | alt. RN | rot. AVL | rot. RN
--------+----------+----------+---------+----------+--------
   1000 |       19 |       11 |      13 |      718 |    1204
  10000 |       31 |       15 |      17 |     7027 |   11970
 100000 |       40 |       19 |      22 |    69951 |  118478
\end{minted}

  \vspace{0.1cm}
  \footnotesize \texttt{go run comparacao.go}
\end{frame}

\begin{frame}{O que a tabela diz}
  \begin{itemize}
    \item as duas eliminam a degeneração, e ficam bem abaixo dos limites
          teóricos (24 e 34 para $n = 100$ mil);
    \item a AVL fica 2 ou 3 níveis mais baixa na ordem aleatória. Cada nível é
          uma comparação a mais por busca;
    \item na inserção, a rubro-negra \tcr{não} faz menos rotações. Na variante
          left-leaning faz mais, e ainda recolore em cerca de três quartos das
          inserções.
  \end{itemize}

  \vspace{0.4cm}
  A vantagem da rubro-negra não está na inserção.
\end{frame}

\begin{frame}{A diferença está na remoção}
  \begin{center}
    \begin{tabular}{lll}
      \toprule
       & \textbf{AVL} & \textbf{Rubro-negra} \\
      \midrule
      Altura no pior caso & $1{,}44\log_2(n+2)$ & $2\log_2(n+1)$ \\
      Rotações por inserção & no máximo 2 & no máximo 2 \\
      Rotações por remoção & até $O(\log n)$ & \tcr{no máximo 3} \\
      Informação por nó & altura ou fator & cor (1 bit) \\
      \bottomrule
    \end{tabular}
  \end{center}

  \vspace{0.4cm}
  Na AVL, o conserto da remoção pode reduzir a altura da subárvore e propagar o
  desbalanceamento para o pai, nível a nível.

  \vspace{0.3cm}
  Muitas remoções: rubro-negra. Quase só buscas: AVL.
\end{frame}

\begin{frame}{Onde a rubro-negra é usada}
  \begin{itemize}
    \item \texttt{std::map} e \texttt{std::set} em C++, nas implementações
          usuais;
    \item \texttt{TreeMap} e \texttt{TreeSet} em Java;
    \item o kernel do Linux, em \texttt{lib/rbtree.c}, usado por vários
          subsistemas.
  \end{itemize}

  \vspace{0.5cm}
  O \texttt{map} de Go fica de fora: ele é uma tabela de dispersão, assunto da
  próxima aula. Em Go, uma estrutura ordenada exige biblioteca externa ou
  implementação própria.

  \vspace{0.3cm}
  \footnotesize A remoção na rubro-negra não é desenvolvida nesta aula e não cai
  na prova além do limite de 3 rotações. Cormen, seção 13.4.
\end{frame}

\begin{frame}{Por que as bibliotecas escolhem a rubro-negra}
  \begin{block}{Mudanças na estrutura em número fixo}
    No máximo 2 rotações por inserção e 3 por remoção. Pesa quando a rotação é
    cara: árvore aumentada (recalcular o valor guardado no nó), estrutura
    persistente (copiar o nó), acesso concorrente (sincronizar).
  \end{block}

  \begin{block}{A biblioteca não conhece o uso}
    Pior caso aceitável na busca e na alteração, sem saber a proporção entre
    as duas.
  \end{block}

  \begin{block}{Tradição}
    STL da HP/SGI, Cormen, e as bibliotecas que vieram depois.
  \end{block}

  \vspace{0.1cm}
  \footnotesize A diferença prática é pequena: com muitas buscas, a AVL costuma
  ganhar (o kernel do Windows usa AVL). A LLRB medida rotaciona mais que a
  rubro-negra clássica das bibliotecas.
\end{frame}

% ==========================================
\section{Árvores B}

\begin{frame}{Quando a árvore binária é a estrutura errada}
  Índice com 100 milhões de chaves, em disco.

  \vspace{0.4cm}
  \begin{itemize}
    \item árvore rubro-negra: altura por volta de 40;
    \item cada nível visitado está em um endereço qualquer do arquivo;
    \item disco rotacional: alguns milissegundos por leitura, mais de 100 ms no
          total. SSD: dezenas de microssegundos, contra nanossegundos da
          memória.
  \end{itemize}

  \vspace{0.4cm}
  E o sistema operacional não lê 16 bytes: lê uma \textbf{página} inteira, de
  4096 bytes. Ler 16 e ler 4096 custa quase o mesmo.

  \vspace{0.3cm}
  A árvore binária usa 16 dos 4096 bytes que a leitura trouxe.
\end{frame}

\begin{frame}{O nó do tamanho da página}
  Chave mais ponteiro ocupando 20 bytes, página de 4096 bytes:

  \vspace{0.2cm}
  \begin{center}
    cerca de \textbf{200 chaves} e \textbf{201 filhos} por nó.
  \end{center}

  \vspace{0.4cm}
\begin{center}
  \small
  \begin{tabular}{rrrrr}
    \toprule
    \textbf{ordem} & \textbf{chaves/nó} & \textbf{1 mil} & \textbf{1 mi} & \textbf{1 bilhão} \\
    \midrule
    5    &    4 & 5 & 9 & 13 \\
    101  &  100 & 2 & 3 &  5 \\
    201  &  200 & 2 & 3 &  4 \\
    1001 & 1000 & 1 & 2 &  3 \\
    \bottomrule
  \end{tabular}
\end{center}

  \vspace{0.3cm}
  Níveis com os nós cheios, ou seja, leituras de disco por busca. Um bilhão de
  chaves em 4 leituras.
\end{frame}

\begin{frame}{Árvore B de ordem $m$}
  \begin{itemize}
    \item cada nó guarda no máximo $m-1$ chaves em ordem crescente, e no máximo
          $m$ filhos;
    \item as chaves do nó dividem as das subárvores em faixas;
    \item todo nó, exceto a raiz, guarda ao menos
          $\lceil m/2 \rceil - 1$ chaves;
    \item \tcr{todas as folhas estão no mesmo nível}.
  \end{itemize}

  \vspace{0.4cm}
  \begin{center}
    \begin{tikzpicture}[-, level distance=1.2cm,
        level 1/.style={sibling distance=2.6cm}]
      \node [bnode] {9}
      child { node [bnode] {3 $\mid$ 6} }
      child { node [bnode] {12 $\mid$ 15 $\mid$ 18} };
    \end{tikzpicture}
  \end{center}

  \vspace{0.2cm}
  O critério é mais forte que o da AVL: não há caminho mais longo que outro.
\end{frame}

\begin{frame}[fragile]{Busca}
\begin{minted}{go}
func Busca(n *NoB, v int) bool {
	if n == nil {
		return false
	}
	i, achou := posicao(n, v)
	if achou {
		return true
	}
	if ehFolha(n) {
		return false
	}
	return Busca(n.Filhos[i], v)
}
\end{minted}

  \vspace{0.2cm}
  Dentro do nó, a busca escolhe entre $m$ caminhos, e não entre 2. Sequencial ou
  binária dá no mesmo: as duas acontecem dentro da página já lida.
\end{frame}

\begin{frame}{Inserção: sempre em folha, com divisão}
  Ordem 5, folha cheia com \texttt{[1 2 4 5]}, chega o 3:

  \vspace{0.3cm}
  \begin{center}
    \begin{tikzpicture}[node distance=0.6cm]
      \node [bnode] (cheio) {1 $\mid$ 2 $\mid$ 3 $\mid$ 4 $\mid$ 5};
      \node [below=0.8cm of cheio] (mid) {};
      \node [bnode, below=1.4cm of cheio] (pai) {3};
      \node [bnode, below left=0.7cm and 0.3cm of pai] (esq) {1 $\mid$ 2};
      \node [bnode, below right=0.7cm and 0.3cm of pai] (dir) {4 $\mid$ 5};
      \draw [-{Latex}] (cheio) -- (pai);
      \draw (pai) -- (esq);
      \draw (pai) -- (dir);
    \end{tikzpicture}
  \end{center}

  \vspace{0.2cm}
  Cinco chaves em um nó que só comporta quatro. A chave do meio \textbf{sobe}
  para o pai, e as outras ficam em dois nós.
\end{frame}

\begin{frame}[fragile]{A divisão em código}
\begin{minted}{go}
func divide(n *NoB) (int, *NoB) {
	meio := len(n.Chaves) / 2
	subiu := n.Chaves[meio]

	direito := &NoB{Chaves: append([]int{},
		n.Chaves[meio+1:]...)}
	if !ehFolha(n) {
		direito.Filhos = append([]*NoB{},
			n.Filhos[meio+1:]...)
		n.Filhos = n.Filhos[:meio+1]
	}
	n.Chaves = n.Chaves[:meio]

	return subiu, direito
}
\end{minted}
\end{frame}

\begin{frame}{A árvore cresce pela raiz}
  A chave que sobe pode estourar o pai, que se divide também. A divisão pode
  propagar até a raiz.

  \vspace{0.4cm}
  \begin{block}{Quando a raiz se divide}
    A árvore ganha um nível. É o único jeito de uma árvore B crescer em altura,
    e é por isso que todas as folhas continuam no mesmo nível.
  \end{block}

  \vspace{0.4cm}
  Uma árvore binária cresce pelas folhas, e por isso um ramo pode ficar mais
  longo que os outros.
\end{frame}

\begin{frame}[fragile]{20 chaves em ordem crescente, ordem 5}
\begin{minted}{text}
$ go run btree.go
inserção do  5: a raiz se dividiu, altura agora é 1
inserção do 17: a raiz se dividiu, altura agora é 2

[9]
    [3 6]
        [1 2]
        [4 5]
        [7 8]
    [12 15 18]
        [10 11]
        [13 14]
        [16 17]
        [19 20]
\end{minted}

  \vspace{0.1cm}
  A mesma entrada que degenera a ABB (altura 19) dá 3 níveis aqui.
\end{frame}

\begin{frame}{B+ e o que os sistemas reais usam}
  Na \textbf{árvore B+}, as chaves ficam todas nas folhas, os nós internos
  guardam apenas separadores, e as folhas são ligadas em lista encadeada.

  \vspace{0.4cm}
  As duas mudanças servem à consulta por faixa: achada a chave inicial, o resto
  do intervalo sai seguindo os ponteiros entre folhas.

  \vspace{0.4cm}
  \begin{itemize}
    \item SQLite: tabelas são árvores B+, índices são árvores B;
    \item PostgreSQL e MySQL/InnoDB: índices em árvore B+;
    \item sistemas de arquivos como Btrfs e XFS.
  \end{itemize}
\end{frame}

\begin{frame}[fragile]{De volta à demonstração}
\begin{minted}{text}
`--SCAN aluno                          69 ms
`--SEARCH aluno USING INDEX idx_grr    < 1 ms
`--SEARCH ... USING COVERING INDEX idx_grr
       (grr>? AND grr<?)
\end{minted}

  \vspace{0.3cm}
  \begin{itemize}
    \item \texttt{SCAN}: lê todas as páginas da tabela, $\Theta(n)$;
    \item \texttt{SEARCH}: desce a árvore B, 3 níveis para 2 milhões de chaves;
    \item o \texttt{BETWEEN} também usa o índice, porque as chaves estão em
          ordem;
    \item os 0,43 s do \texttt{CREATE INDEX} são 2 milhões de inserções em
          árvore B. Indexar acelera a consulta e custa na escrita.
  \end{itemize}
\end{frame}

% ==========================================
\section{Tries}

\begin{frame}{Chaves que têm estrutura interna}
  Até aqui, toda estrutura compara chaves inteiras. Com chaves de texto, cada
  comparação percorre a palavra até a primeira letra diferente, e a árvore de
  altura $\log n$ faz $\log n$ dessas comparações.

  \vspace{0.5cm}
  A \textbf{trie} (de re\textbf{trie}val) usa uma letra da chave por nível para
  escolher o filho.

  \vspace{0.3cm}
  A chave não fica guardada em nó nenhum: ela é o \textbf{caminho} da raiz até o
  nó, e um marcador diz quais caminhos são palavras inteiras.
\end{frame}

\begin{frame}{Uma trie com cinco palavras}
  \begin{columns}
    \begin{column}{0.55\textwidth}
      \centering
      \begin{tikzpicture}[-, level distance=0.75cm,
          level 1/.style={sibling distance=2.4cm},
          level 2/.style={sibling distance=1.6cm},
          level 3/.style={sibling distance=1.2cm}]
        \node [tnode] {}
        child { node [tnode] {c}
            child { node [tnode] {a}
                child { node [tnode] {s}
                    child { node [tfim] {a}
                        child { node [tfim] {l} }
                      }
                  }
                child { node [tnode] {r}
                    child { node [tnode] {t}
                        child { node [tfim] {a}
                            child { node [tfim] {z} }
                          }
                      }
                  }
              }
          }
        child { node [tnode] {d}
            child { node [tnode] {a}
                child { node [tnode] {d}
                    child { node [tfim] {o} }
                  }
              }
          };
      \end{tikzpicture}
    \end{column}
    \begin{column}{0.4\textwidth}
      \texttt{casa}, \texttt{casal}, \texttt{carta}, \texttt{cartaz},
      \texttt{dado}.

      \vspace{0.3cm}
      Nó \tcg{verde} marca fim de palavra.

      \vspace{0.3cm}
      13 nós além da raiz, para 24 letras. Os prefixos comuns são guardados uma
      vez só.
    \end{column}
  \end{columns}
\end{frame}

\begin{frame}[fragile]{Nó e inserção}
\begin{minted}{go}
type No struct {
	Filhos       map[rune]*No
	FimDePalavra bool
}

func Insere(raiz *No, palavra string) {
	atual := raiz
	for _, letra := range palavra {
		proximo, existe := atual.Filhos[letra]
		if !existe {
			proximo = NovoNo()
			atual.Filhos[letra] = proximo
		}
		atual = proximo
	}
	atual.FimDePalavra = true
}
\end{minted}
\end{frame}

\begin{frame}{O custo não depende do número de chaves}
  \begin{center}
    \small
    \begin{tabular}{lll}
      \toprule
      \textbf{Operação} & \textbf{Trie} & \textbf{ABB balanceada} \\
      \midrule
      Busca de palavra de $m$ letras & $O(m)$ & $O(m \log n)$ \\
      Consulta por prefixo & $O(m + k)$ & percorre a faixa \\
      Espaço & um nó por letra & um nó por chave \\
      \bottomrule
    \end{tabular}
  \end{center}

  \vspace{0.4cm}
  Mil ou dez milhões de palavras: a busca custa o mesmo, porque o caminho
  percorrido é a própria palavra.

  \vspace{0.3cm}
  Em troca, a trie ocupa mais espaço. O nó pode ser um \texttt{map} (pouco
  espaço, poucos filhos) ou um vetor com uma posição por letra do alfabeto
  (acesso constante, posições desperdiçadas).
\end{frame}

\begin{frame}[fragile]{Consulta por prefixo}
\begin{minted}{text}
$ go run trie.go
17 palavras, 86 letras no total, 37 nós na trie

Contem("cas") = false, EhPrefixo("cas") = true

ComPrefixo("car") = carro, carta, cartaz, cartão
ComPrefixo("cas") = casa, casaco, casal, caso, cassino
ComPrefixo("z") = nenhuma palavra
\end{minted}

  \vspace{0.2cm}
  Desce até o nó do prefixo e coleta a subárvore. Nenhuma palavra fora da
  resposta é visitada.

  \vspace{0.2cm}
  Autocompletar, corretor ortográfico e roteamento IP por prefixo mais longo.
  Tabela de dispersão não faz isso: o hash destrói a vizinhança entre chaves
  parecidas.
\end{frame}

% ==========================================
\section{Fechamento}

\begin{frame}{Escolher entre as estruturas}
  \begin{center}
    \small
    \begin{tabular}{llp{4.2cm}}
      \toprule
      \textbf{Estrutura} & \textbf{Busca} & \textbf{Escolha quando} \\
      \midrule
      ABB sem balanceamento & $O(h)$ & nunca, em produção \\
      AVL & $\Theta(\log n)$ & busca domina, alterações raras \\
      Rubro-negra & $\Theta(\log n)$ & uso geral, muitas alterações \\
      Árvore B / B+ & $\Theta(\log_m n)$ & a estrutura vive em disco \\
      Trie & $O(m)$ & chaves são cadeias, consulta por prefixo \\
      \bottomrule
    \end{tabular}
  \end{center}

  \vspace{0.4cm}
  Nas três primeiras linhas, o critério é a proporção entre buscas e alterações.
  Na quarta, o meio em que a estrutura vive. Na quinta, o tipo da chave.
\end{frame}

\begin{frame}{Para a próxima aula}
  \begin{itemize}
    \item exercícios 1 a 12 do roteiro da aula, sem entrega e sem nota. O
          conteúdo cai na Prova 2;
    \item rodar \texttt{rn.go}, \texttt{comparacao.go}, \texttt{btree.go} e
          \texttt{trie.go}, alterando as sequências de inserção;
    \item leitura: Sedgewick, seções 3.3 e 5.2; Cormen, capítulos 13 e 18.
  \end{itemize}

  \vspace{0.5cm}
  Próxima aula: tabelas de dispersão. A estrutura que abre mão da ordem entre as
  chaves para buscar em tempo constante.
\end{frame}

\end{document}
